\documentclass{beamer}

\usetheme{Berlin}
\usecolortheme{Orchid}
\useoutertheme{miniframes}
\setbeamertemplate{navigation symbols}{}

\usepackage[utf8]{inputenc}
\usepackage[T1]{fontenc}
\usepackage{amsmath}
\usepackage{amssymb}
\usepackage{booktabs}
\usepackage{graphicx}
\graphicspath{{../images/}}

\usepackage{xcolor}
\usepackage{listings}
\usepackage{hyperref}
\usepackage{tikz}
\usetikzlibrary{arrows.meta,positioning,shapes.geometric,calc}

\lstset{
  language=Java,
  basicstyle=\ttfamily\small,
  keywordstyle=\color{blue},
  commentstyle=\color{gray},
  stringstyle=\color{teal},
  showstringspaces=false,
  columns=fullflexible,
  breaklines=true
}

% Per-slide source line (references shown on the slide itself).
\newcommand{\slidesources}[1]{%
  \vfill
  {\tiny\color{gray}\raggedright\sloppy\parbox{\linewidth}{Sources: #1}\par}%
}

% Unit-wise source strings for this subject (used by slides/notes).
% Path is relative to the lecture `latex/` folder (the compilation CWD).
% OOPS with Java (CSEG1044) — unit-wise sources
%
% These are short, student-facing reference strings intended to be used with:
%   - \slidesources{...}  (in beamer slides)
%   - \notesources{...}   (in lecture notes)
%
% Style inspiration: other subjects in My-Subjects/ use semicolon-separated sources
% and include an "accessed" date for web documentation.

\newcommand{\OOPSUnitISources}{Schildt, \textit{Java: A Beginner's Guide} (10e), McGraw-Hill (Java fundamentals + OOP basics).; Oracle Java Tutorials: Getting Started + Language Basics + Classes and Objects (accessed \today).; Java SE API docs (e.g., \texttt{String}, \texttt{Scanner}) (accessed \today).}

\newcommand{\OOPSUnitIISources}{Schildt, \textit{Java: The Complete Reference} (12e), McGraw-Hill (classes, inheritance, packages, interfaces).; Bloch, \textit{Effective Java} (3e), Addison-Wesley, 2018 (design choices; interfaces vs abstract classes).; Oracle Java Tutorials: Classes and Objects + Interfaces + Packages (accessed \today).; Java SE API docs (e.g., \texttt{Comparable}, \texttt{Runnable}) (accessed \today).}

\newcommand{\OOPSUnitIIISources}{Schildt, \textit{Java: The Complete Reference} (12e) (nested/inner classes, exceptions, threads, I/O).; Oracle Java Tutorials: Nested Classes + Exceptions + Concurrency + Basic I/O (accessed \today).; Java SE API docs (e.g., \texttt{Thread}, \texttt{AutoCloseable}) (accessed \today).}

\newcommand{\OOPSUnitIVSources}{Schildt, \textit{Java: The Complete Reference} (12e) (generics, lambdas, Swing, JDBC).; Bloch, \textit{Effective Java} (3e) (generics + lambdas).; Oracle Java Tutorials: Generics + Lambda Expressions + Swing + JDBC Basics (accessed \today).; Java SE API docs (e.g., \texttt{java.util.function}, \texttt{javax.swing}, \texttt{java.sql}) (accessed \today).}

\newcommand{\OOPSUnitVSources}{Schildt, \textit{Java: The Complete Reference} (12e) (Collections Framework + wrapper classes).; Oracle Java docs: Collections Framework overview (accessed \today).; Java SE API docs (\texttt{java.util} collections; wrapper classes in \texttt{java.lang}) (accessed \today).}

\newcommand{\OOPSUnitVISources}{Git SCM: \textit{Pro Git} (2e) (version control workflow), accessed \today.; GitHub Docs (repositories, issues, pull requests), accessed \today.; Oracle Java Tutorials: Swing + JDBC (accessed \today).; JUnit 5 User Guide (accessed \today).; Apache Maven Project: Getting Started (accessed \today).; Gradle User Manual: Getting Started (accessed \today).; UML class diagrams: UML 2.x overview/spec (accessed \today).}

\newcommand{\OOPSMiscWhyInterfacesSources}{Schildt, \textit{Java: The Complete Reference} (12e) (interfaces).; Bloch, \textit{Effective Java} (3e) (prefer interfaces where appropriate).; Oracle Java Tutorials: Interfaces (accessed \today).; Java SE API docs (e.g., \texttt{Comparable}, \texttt{Runnable}) (accessed \today).}



\title[OOPS with Java]{Object Oriented Programming (CSEG1044)}
\subtitle{Unit 03 -- Lecture 02: Exception Handling (Core)}
\author{Tofik Ali}
\institute{School of Computer Science, UPES Dehradun}
\date{\today}

\begin{document}

\begin{frame}[fragile]
  \titlepage
  \vspace{0.5em}
  \slidesources{\OOPSUnitIIISources}
\end{frame}

\begin{frame}{Quick Links}
  \centering
  \hyperlink{sec:overview}{\beamerbutton{Overview}}\hspace{0.6em}
  \hyperlink{sec:concepts}{\beamerbutton{Key Topics}}\hspace{0.6em}
  \hyperlink{sec:demo}{\beamerbutton{Demo}}\hspace{0.6em}
  \hyperlink{sec:summary}{\beamerbutton{Summary}}
  \slidesources{\OOPSUnitIIISources}
\end{frame}

\begin{frame}{Agenda}
  \tableofcontents
  \slidesources{\OOPSUnitIIISources}
\end{frame}

\section{Overview}
\label{sec:overview}

\begin{frame}{Learning Outcomes}
  \begin{itemize}[<+->]
    \item Distinguish exception vs error; checked vs unchecked.
    \item Write \texttt{try/catch/finally} and use \texttt{throw} vs \texttt{throws}.
    \item Design a small custom exception for validation.
  \end{itemize}
  \slidesources{\OOPSUnitIIISources}
\end{frame}

\section{Key Topics}
\label{sec:concepts}

\begin{frame}{Unit 03: Nested Classes, Exceptions, Multithreading \& IO Streams}
  \begin{itemize}[<+->]
    \item Lecture focus: Exception Handling (Core)
  \end{itemize}
  \slidesources{\OOPSUnitIIISources}
\end{frame}

\begin{frame}{Today: Roadmap}
  \begin{itemize}[<+->]
    \item Exceptions vs errors; checked vs unchecked
    \item `try/catch/finally`, multiple catch; `throw` vs `throws`
    \item Exception handlers + basic custom exception design for validation
  \end{itemize}
  \slidesources{\OOPSUnitIIISources}
\end{frame}

\begin{frame}{Exception vs Error}
  \begin{itemize}[<+->]
    \item Exceptions: conditions you can handle (invalid input, missing file).
    \item Errors: serious JVM/system problems (usually not recoverable).
  \end{itemize}
  \slidesources{\OOPSUnitIIISources}
\end{frame}

\begin{frame}{Checked vs Unchecked}
  \begin{itemize}[<+->]
    \item Checked: must be caught or declared (\texttt{throws}) to compile.
    \item Unchecked: \texttt{RuntimeException} (compiler does not force handling).
  \end{itemize}
  \slidesources{\OOPSUnitIIISources}
\end{frame}

\begin{frame}{\texttt{try/catch/finally} (pattern)}
  \begin{itemize}[<+->]
    \item Put risky code in \texttt{try}.
    \item Catch specific exceptions first.
    \item Use \texttt{finally} for cleanup (runs even if exception occurs).
  \end{itemize}
  \slidesources{\OOPSUnitIIISources}
\end{frame}

\begin{frame}{\texttt{throw} vs \texttt{throws}}
  \begin{itemize}[<+->]
    \item \texttt{throw}: create and throw an exception now.
    \item \texttt{throws}: declare that a method may throw an exception.
  \end{itemize}
  \slidesources{\OOPSUnitIIISources}
\end{frame}

\begin{frame}[fragile]{Example: \texttt{throw} (happens now)}
\begin{lstlisting}
static int divide(int a, int b) {
    if (b == 0) {
        throw new IllegalArgumentException("b must not be 0");
    }
    return a / b;
}
\end{lstlisting}
  \begin{itemize}[<+->]
    \item Thrown at runtime when \texttt{b == 0}.
    \item Execution jumps to a matching \texttt{catch}.
  \end{itemize}
  \slidesources{\OOPSUnitIIISources}
\end{frame}

\begin{frame}[fragile]{Example: \texttt{throws} (declared in signature)}
\begin{lstlisting}
static void pauseOneSecond() throws InterruptedException {
    Thread.sleep(1000); // checked exception
}
\end{lstlisting}
  \begin{itemize}[<+->]
    \item \texttt{throws} tells the caller: catch or declare.
    \item Checked exceptions are enforced by the compiler.
  \end{itemize}
  \slidesources{\OOPSUnitIIISources}
\end{frame}

\begin{frame}[fragile]{Caller must handle (try/catch) or declare (throws)}
\begin{lstlisting}
public static void main(String[] args) {
    try {
        pauseOneSecond();
    } catch (InterruptedException ex) {
        System.out.println("Interrupted");
    }
}
\end{lstlisting}
  \slidesources{\OOPSUnitIIISources}
\end{frame}


\section{Demo / Example}
\label{sec:demo}

\begin{frame}[fragile]{Full example (1/2): \texttt{try} block}
\begin{lstlisting}
import java.util.Scanner;

public class DivideDemo {
    public static void main(String[] args) {
        Scanner sc = new Scanner(System.in);

        try {
            System.out.print("Enter a: ");
            int a = Integer.parseInt(sc.nextLine());
            System.out.print("Enter b: ");
            int b = Integer.parseInt(sc.nextLine());

            int q = a / b;
            System.out.println("a/b = " + q);
\end{lstlisting}
  \slidesources{\OOPSUnitIIISources}
\end{frame}

\begin{frame}[fragile]{Full example (2/2): \texttt{catch} + \texttt{finally}}
\begin{lstlisting}
        } catch (NumberFormatException ex) {
            System.out.println("Please enter valid integers.");
        } catch (ArithmeticException ex) {
            System.out.println("Cannot divide by zero.");
        } finally {
            sc.close();
            System.out.println("Cleanup done (finally).");
        }
    }
}
\end{lstlisting}
  \slidesources{\OOPSUnitIIISources}
\end{frame}

\begin{frame}{How the example behaves}
  \begin{itemize}[<+->]
    \item Valid input: prints quotient, then runs \texttt{finally}.
    \item Invalid integer input: \texttt{NumberFormatException} catch runs, then \texttt{finally}.
    \item Divide by zero: \texttt{ArithmeticException} catch runs, then \texttt{finally}.
  \end{itemize}
  \slidesources{\OOPSUnitIIISources}
\end{frame}

\begin{frame}{Execution sequence (try/catch/finally)}
  \begin{enumerate}[<+->]
    \item Enter \texttt{try} and execute statements top-to-bottom.
    \item If no exception: skip all \texttt{catch} blocks.
    \item If an exception occurs: jump to the \textbf{first matching} \texttt{catch}.
    \item After \texttt{try} or \texttt{catch}: always execute \texttt{finally} (cleanup).
    \item Continue with the next statement after the block (or propagate if not handled).
  \end{enumerate}
  \slidesources{\OOPSUnitIIISources}
\end{frame}

\begin{frame}[fragile]{Example: custom checked exception}
\begin{lstlisting}
class InsufficientFundsException extends Exception {
    InsufficientFundsException(String msg) { super(msg); }
}

class BankAccount {
    private double balance;
    BankAccount(double opening) { balance = opening; }
    void withdraw(double amount) throws InsufficientFundsException {
        if (amount > balance) throw new InsufficientFundsException("need " + amount);
        balance -= amount;
    }
}
\end{lstlisting}
  \slidesources{\OOPSUnitIIISources}
\end{frame}

\begin{frame}[fragile]{Handling it}
\begin{lstlisting}
public class Demo {
    public static void main(String[] args) {
        BankAccount a = new BankAccount(100);
        try {
            a.withdraw(150);
        } catch (InsufficientFundsException e) {
            System.out.println("Withdraw failed: " + e.getMessage());
        }
    }
}
\end{lstlisting}
  \slidesources{\OOPSUnitIIISources}
\end{frame}

\begin{frame}{Mini Exercises (2--3 mins)}
  \begin{enumerate}[<+->]
    \item Give one checked and one unchecked exception example.
    \item When does \texttt{finally} run?
    \item Difference between \texttt{throw} and \texttt{throws}?
  \end{enumerate}
  \slidesources{\OOPSUnitIIISources}
\end{frame}


\section{Summary}
\label{sec:summary}

\begin{frame}{Key Takeaways}
  \begin{itemize}[<+->]
    \item Exceptions enable safe error handling without crashing.
    \item Catch specific exceptions; use \texttt{finally} for cleanup.
    \item Custom exceptions make domain errors clear.
  \end{itemize}
  \slidesources{\OOPSUnitIIISources}
\end{frame}

\begin{frame}{Checkpoint Questions}
  \begin{enumerate}[<+->]
    \item Why catch specific exceptions (not \texttt{Exception})?
    \item What does ``checked exception'' mean?
    \item When do we use \texttt{throws} in a method signature?
  \end{enumerate}
  \slidesources{\OOPSUnitIIISources}
\end{frame}

\begin{frame}{Next Lecture Preview}
  \textbf{Next:} Multithreading + Synchronization
  \vspace{0.6em}
  \begin{itemize}[<+->]
    \item We will create threads using \texttt{Thread} and \texttt{Runnable} and see race conditions.
  \end{itemize}
  \slidesources{\OOPSUnitIIISources}
\end{frame}

\end{document}
